\subsection{Crossovers between regimes}

Regimes 1--4 each describe a limit in which one or two terms of antibody
conservation dominate. Real experiments need not sit deep inside any one
limit, and the boundaries between regimes are exactly where the
treatment and control conditions are most likely to behave differently.
In this section we first prove a result that holds in \emph{every}
regime, and then work out the three boundary regions explicitly:
\begin{itemize}
    \item the \textbf{equivalence crossover} ($Y\approx S$), between
    regimes 1 and 3;
    \item the \textbf{off-target half-saturation crossover}
    ($Y_f\sim K_{d2}$), between regimes 3 and 4;
    \item the \textbf{weak-binding crossover} ($K_{d1}\sim S$), between
    regimes 1 and 2.
\end{itemize}
As before, $S\in\{S_c,S_t\}$ with $S_c=A+B$ and $S_t=A'+B$, and
$r=K_{d1}/K_{d2}$. We abbreviate the off-target pool as $M=C+D$, which is
the same in both conditions.

\subsubsection{A general principle: less target means more free antibody}

\textbf{Step 1: $\rho$ depends only on $Y_f$, and increases with it.}
By definition,
$$
\rho=\frac{Y_f+K_{d1}}{Y_f+K_{d2}},
$$
which involves no concentration other than $Y_f$. Treating $\rho$ as a
function of $Y_f$ and differentiating with the quotient rule,
$$
\frac{d\rho}{dY_f}
=\frac{1\cdot(Y_f+K_{d2})-(Y_f+K_{d1})\cdot1}{(Y_f+K_{d2})^2}
=\frac{K_{d2}-K_{d1}}{(Y_f+K_{d2})^2}.
$$
Since $K_{d1}<K_{d2}$, the numerator is positive, and the denominator is a
square, so
\begin{equation}
\frac{d\rho}{dY_f}=\frac{K_{d2}-K_{d1}}{(Y_f+K_{d2})^2}>0.
\label{eq:x-drho}
\end{equation}
Thus, \textbf{whichever condition has more free antibody has the larger
$\rho$}. Intuitively, free antibody is the ``pressure'' that pushes
antibody onto epitopes. Target epitopes, with their small $K_{d1}$, are
already mostly occupied at low pressure, so raising the pressure helps
the weak off-target epitopes more than the target epitopes.

\textbf{Step 2: $Y_f$ is set implicitly by antibody conservation.}
Substituting $\theta_1=Y_f/(Y_f+K_{d1})$ and $\theta_2=Y_f/(Y_f+K_{d2})$
into antibody conservation $Y=Y_f+S\theta_1+M\theta_2$ gives
\begin{equation}
Y=F(Y_f;S),
\qquad
F(y;S)=y+\frac{Sy}{y+K_{d1}}+\frac{My}{y+K_{d2}}.
\label{eq:x-F}
\end{equation}
Here $F(y;S)$ is the total antibody that would be needed to produce a free
antibody concentration $y$ when the total target is $S$: $y$ free, plus
$Sy/(y+K_{d1})$ bound to target, plus $My/(y+K_{d2})$ bound off-target.
We cannot in general solve \eqref{eq:x-F} for $Y_f$ in closed form (after
clearing denominators it is a cubic), but we can learn a lot from how $F$
depends on its two arguments.

\emph{$F$ is increasing in $y$.} Each fraction has the form
$y/(y+K)$, whose derivative by the quotient rule is
$$
\frac{d}{dy}\left(\frac{y}{y+K}\right)
=\frac{1\cdot(y+K)-y\cdot1}{(y+K)^2}
=\frac{K}{(y+K)^2}.
$$
Applying this with $K=K_{d1}$ and $K=K_{d2}$,
\begin{equation}
\frac{\partial F}{\partial y}
=1+\frac{SK_{d1}}{(y+K_{d1})^2}+\frac{MK_{d2}}{(y+K_{d2})^2}
\ge1>0.
\label{eq:x-dFdy}
\end{equation}
(The symbol $\partial$ means we differentiate with respect to $y$ while
holding $S$ fixed.) Also $F(0;S)=0$ and $F(y;S)\ge y$, so $F$ climbs
continuously from $0$ to arbitrarily large values. By the intermediate
value theorem, for any $Y>0$ there is a solution $Y_f$ of
$F(Y_f;S)=Y$, and because $F$ is strictly increasing in $y$, it is the only
one. Moreover, $Y_f>0$, since $F(0;S)=0<Y$.

\emph{$F$ is increasing in $S$.} $S$ appears only in the middle term, so
$$
\frac{\partial F}{\partial S}=\frac{y}{y+K_{d1}}=\theta_1>0
\quad\text{for } y>0.
$$
More antibody is needed to reach a given free concentration when there
is more target to soak it up.

\textbf{Step 3: smaller $S$ gives larger $Y_f$.} Let $Y_f$ and $Y_f'$ be
the solutions for the control ($S_c$) and treatment ($S_t$), with the same
$Y$, and suppose $A'<A$, so $S_t<S_c$. Evaluate the treatment's function
at the control's free antibody. Since only the middle term depends on
$S$,
$$
F(Y_f;S_t)
=F(Y_f;S_c)-\frac{(S_c-S_t)\,Y_f}{Y_f+K_{d1}}
=Y-(S_c-S_t)\,\theta_1,
$$
where we substituted $F(Y_f;S_c)=Y$ from \eqref{eq:x-F}. The subtracted
term is positive (because $S_c>S_t$ and $\theta_1>0$), so
$$
F(Y_f;S_t)<Y=F(Y_f';S_t).
$$
Since $F(\,\cdot\,;S_t)$ is increasing \eqref{eq:x-dFdy}, a smaller output
requires a smaller input, so $Y_f<Y_f'$. In words: removing target
epitopes releases the antibody that would have bound them, and the
system re-equilibrates at a higher free-antibody concentration.

\textbf{Step 4: result.} Combining Step 3 with \eqref{eq:x-drho},
\begin{equation}
A'<A
\quad\Longrightarrow\quad
Y_f'>Y_f
\quad\Longrightarrow\quad
\rho'>\rho.
\label{eq:x-mono}
\end{equation}
This holds in \emph{every} regime, without any approximation. The
regime-specific formulas only tell us \emph{how much} larger $\rho'$ is:
in the deep limits of regimes 2 and 4 the difference is negligible
($\rho'\approx\rho$), whereas near the crossovers below it can be large.

\textbf{How large can the difference be?} The same calculation can be
phrased with implicit differentiation. Differentiating both sides of
$Y=F(Y_f(S);S)$ with respect to $S$ (with $Y$ fixed) using the chain rule,
$$
0=\frac{\partial F}{\partial y}\frac{dY_f}{dS}+\frac{\partial F}{\partial S},
\qquad\text{so}\qquad
\frac{dY_f}{dS}
=-\frac{\partial F/\partial S}{\partial F/\partial y}
=-\frac{\theta_1}{\partial F/\partial y}.
$$
Since $0<\theta_1\le1$ and $\partial F/\partial y\ge1$ by
\eqref{eq:x-dFdy}, we get $-1\le dY_f/dS<0$. By the mean value theorem,
there is some $S^*$ between $S_t$ and $S_c$ with
$$
Y_f'-Y_f
=-\frac{dY_f}{dS}\bigg|_{S^*}\,(S_c-S_t)
\le S_c-S_t
=A-A'.
$$
That is, removing $A-A'$ target epitopes can raise the free antibody by at
most $A-A'$ (the released antibody can only be shared out among solution,
the remaining target, and off-target sites). Applying the mean value
theorem again, now to $\rho$ as a function of $Y_f$, and bounding the
derivative \eqref{eq:x-drho} by its largest value (at $Y_f=0$),
\begin{equation}
\rho'-\rho
\le\frac{K_{d2}-K_{d1}}{K_{d2}^2}\,(Y_f'-Y_f)
\le\frac{(1-r)(A-A')}{K_{d2}}
<\frac{A-A'}{K_{d2}}.
\label{eq:x-bound}
\end{equation}
So the \emph{absolute} change in $\rho$ is always modest unless
$A-A'$ is comparable to $K_{d2}$. The \emph{relative} change
$\rho'/\rho$ can nevertheless be large when $\rho$ itself is small
($\rho\sim r$) and $Y_f$ is very sensitive to $S$, which is exactly the
situation near equivalence.

\subsubsection{Equivalence crossover (between regimes 1 and 3)}

Regime 1 required $\sqrt{K_{d1}Y}\ll S-Y$, and regime 3 requires the
antibody to exceed the target. Between them lies a window around
$Y\approx S$ where antibody and target \emph{titrate} each other. Define
the excess target
$$
\delta=S-Y,
$$
which is positive on the regime-1 side and negative on the regime-3 side.

\textbf{Assumptions.}
\begin{enumerate}
    \item Off-target-bound antibody is negligible compared with free
    antibody: $M\theta_2\ll Y_f$. Since
    $M\theta_2=MY_f/(Y_f+K_{d2})\le MY_f/K_{d2}$, a sufficient condition
    is $M\ll K_{d2}$. (We relax this in a remark below.)
    \item Target binding is tight compared with the amount of target:
    $K_{d1}\ll S$.
    \item We are near equivalence: $|\delta|\ll S$.
    \item Off-target epitopes are far from half-occupied: $Y_f\ll K_{d2}$.
\end{enumerate}

\textbf{Step 1: an exact quadratic.} Dropping $M\theta_2$ (assumption 1)
from \eqref{eq:r1-1}, the free target is
$$
S(1-\theta_1)=(S-Y)+Y_f=\delta+Y_f,
$$
and solving antibody conservation for the bound target gives
$S\theta_1=Y-Y_f$. The mass-action relation for target,
$K_{d1}=Y_f(A_f+B_f)/(Y_A+Y_B)$, has free target $A_f+B_f=S(1-\theta_1)$
and bound target $Y_A+Y_B=S\theta_1$. Substituting both and multiplying
through by $Y-Y_f$,
$$
K_{d1}(Y-Y_f)=Y_f(\delta+Y_f).
$$
Expanding and collecting everything on one side,
\begin{equation}
Y_f^2+(\delta+K_{d1})Y_f-K_{d1}Y=0.
\label{eq:x-quad-exact}
\end{equation}
Given assumption 1, this is exact; we will reuse it for the weak-binding
crossover.

\textbf{Step 2: simplify near equivalence.} Two terms in
\eqref{eq:x-quad-exact} are small corrections near equivalence. First,
substituting $Y=S-\delta$,
$$
K_{d1}Y=K_{d1}S\left(1-\frac{\delta}{S}\right)\approx K_{d1}S
$$
by assumption 3. Second, comparing $K_{d1}Y_f$ with $K_{d1}Y$, their ratio
is $Y_f/Y$, which is $\ll1$ provided $Y_f\ll Y\approx S$ (we verify this in
Step 4). Dropping $K_{d1}Y_f$ and replacing $K_{d1}Y$ by $K_{d1}S$,
\begin{equation}
Y_f\,(Y_f+\delta)\approx K_{d1}S.
\label{eq:x-quad}
\end{equation}
This has a direct interpretation. $Y_f+\delta$ is the free target, and
nearly all of the $\approx S$ target epitopes are bound, so
\eqref{eq:x-quad} is just mass action:
(free antibody)$\times$(free target) $=K_{d1}\times$(bound target).

\textbf{Step 3: solve the quadratic.} Writing \eqref{eq:x-quad} as
$Y_f^2+\delta Y_f-K_{d1}S=0$, the quadratic formula gives
$$
Y_f=\frac{-\delta\pm\sqrt{\delta^2+4K_{d1}S}}{2}.
$$
The product of the two roots equals the constant term, $-K_{d1}S<0$, so one
root is positive and one is negative. Free antibody must be positive, so
we take the $+$ sign. Multiplying the numerator and denominator by the
conjugate $\delta+\sqrt{\delta^2+4K_{d1}S}$ and using
$(\sqrt{\delta^2+4K_{d1}S})^2-\delta^2=4K_{d1}S$ gives an equivalent form
that is convenient when $\delta>0$:
\begin{equation}
Y_f
=\frac{-\delta+\sqrt{\delta^2+4K_{d1}S}}{2}
=\frac{2K_{d1}S}{\delta+\sqrt{\delta^2+4K_{d1}S}}.
\label{eq:x-yf}
\end{equation}
The natural scale is $\sqrt{K_{d1}S}$. Writing
$\delta=x\sqrt{K_{d1}S}$ and dividing \eqref{eq:x-yf} by $\sqrt{K_{d1}S}$,
\begin{equation}
\frac{Y_f}{\sqrt{K_{d1}S}}
=h(x)
=\frac{-x+\sqrt{x^2+4}}{2}
=\frac{2}{x+\sqrt{x^2+4}}.
\label{eq:x-h}
\end{equation}
Comparing the two forms of $h$, $h(-x)=(x+\sqrt{x^2+4})/2=1/h(x)$: excess
target and excess antibody of the same size shift $Y_f$ by the same factor
in opposite directions.

\textbf{Step 4: the three limits.}
\begin{itemize}
    \item \emph{Excess target, $\delta\gg\sqrt{K_{d1}S}$ (regime 1).} Then
    $4K_{d1}S\ll\delta^2$, so $\sqrt{\delta^2+4K_{d1}S}\approx\delta$.
    Substituting into the second form of \eqref{eq:x-yf},
    $$
    Y_f\approx\frac{2K_{d1}S}{\delta+\delta}=\frac{K_{d1}S}{\delta}
    =\frac{K_{d1}S}{S-Y}.
    $$
    Since $S\approx Y$ (assumption 3), this matches the regime-1 estimate
    $K_{d1}Y/(S-Y)$ of \eqref{eq:r1-4}. (The exact quadratic
    \eqref{eq:x-quad-exact} gives $K_{d1}Y/(S-Y+K_{d1})$ in this limit,
    matching even more closely.)
    \item \emph{Exact equivalence, $\delta=0$.} Substituting
    $\delta=0$ into \eqref{eq:x-yf},
    $$
    Y_f=\frac{\sqrt{4K_{d1}S}}{2}=\sqrt{K_{d1}S},
    $$
    the geometric mean of $K_{d1}$ and $S$, as anticipated in the
    discussion of regime 1's assumption 2.
    \item \emph{Excess antibody, $-\delta\gg\sqrt{K_{d1}S}$ (regime 3).}
    Now $\delta=-|\delta|$ and again $\sqrt{\delta^2+4K_{d1}S}\approx|\delta|$.
    Substituting into the first form of \eqref{eq:x-yf},
    $$
    Y_f\approx\frac{|\delta|+|\delta|}{2}=|\delta|=Y-S.
    $$
    The target is saturated, and the excess antibody is free (it has
    nowhere else to go, since off-target binding is negligible).
\end{itemize}
We can now verify the claim $Y_f\ll S$ used in Step 2. Using
$\sqrt{a+b}\le\sqrt a+\sqrt b$ for $a,b\ge0$ in the first form of
\eqref{eq:x-yf},
$$
Y_f\le\frac{-\delta+|\delta|+2\sqrt{K_{d1}S}}{2}\le|\delta|+\sqrt{K_{d1}S}.
$$
Here $|\delta|\ll S$ by assumption 3, and $\sqrt{K_{d1}S}\ll S$ because,
dividing by $\sqrt S$, it is equivalent to $\sqrt{K_{d1}}\ll\sqrt S$, i.e.,
assumption 2. So $Y_f\ll S$.

\textbf{Step 5: $\rho$.} By assumption 4, $Y_f+K_{d2}\approx K_{d2}$, so
\begin{equation}
\rho
=\frac{Y_f+K_{d1}}{Y_f+K_{d2}}
\approx\frac{Y_f+K_{d1}}{K_{d2}}
=r\left(1+\frac{Y_f}{K_{d1}}\right),
\label{eq:x-rho-eq}
\end{equation}
where the last step factors out $K_{d1}/K_{d2}=r$. Substituting the three
limits of Step 4:
$$
\rho\approx
\begin{cases}
r\left(1+\dfrac{S}{S-Y}\right)\approx r\dfrac{S}{S-Y}, & S-Y\gg\sqrt{K_{d1}S}
\quad\text{(regime 1)},\\[2ex]
r\left(1+\sqrt{\dfrac{S}{K_{d1}}}\right)\approx\dfrac{\sqrt{K_{d1}S}}{K_{d2}}, & Y=S,\\[2ex]
r+\dfrac{Y-S}{K_{d2}}, & Y-S\gg\sqrt{K_{d1}S}
\quad\text{(regime 3, }M\ll K_{d2}\text{)}.
\end{cases}
$$
(In the first line, $S/(S-Y)\gg1$ near equivalence, so the $1$ is
negligible; in the second, $\sqrt{S/K_{d1}}\gg1$ by assumption 2.)
Through the window, $\rho$ climbs from about $r\,S/(S-Y)$ to about
$(Y-S)/K_{d2}$, and the climb happens over a range of $Y$ of width a few
times $\sqrt{K_{d1}S}$, i.e., a \emph{relative} width of order
$\sqrt{K_{d1}S}/S=\sqrt{K_{d1}/S}\ll1$. For example, by \eqref{eq:x-h},
moving from $x=2$ to $x=-2$ multiplies $Y_f$ by
$h(-2)/h(2)=1/h(2)^2=(1+\sqrt2)^2\approx5.8$; moving from $x=4$ to $x=-4$
multiplies it by $\bigl((4+\sqrt{20})/2\bigr)^2\approx18$.

\textbf{$\rho'$ versus $\rho$.} The treatment has
$\delta_t=S_t-Y=\delta_c-(A-A')<\delta_c$. In \eqref{eq:x-yf} (second
form), the denominator $\delta+\sqrt{\delta^2+4K_{d1}S}$ is increasing in
$\delta$, since its derivative is
$1+\delta/\sqrt{\delta^2+4K_{d1}S}>0$ (because
$|\delta|<\sqrt{\delta^2+4K_{d1}S}$). So decreasing $\delta$ increases
$Y_f$, and by \eqref{eq:x-rho-eq}, $\rho'>\rho$, as guaranteed in general by
\eqref{eq:x-mono}. Because $\rho\propto Y_f+K_{d1}$ here, the ratio is
$$
\frac{\rho'}{\rho}\approx\frac{Y_f'+K_{d1}}{Y_f+K_{d1}},
$$
which is a ratio of free-antibody concentrations. If the drop $A-A'$ is
large enough to carry the treatment across the window while the control
stays on the regime-1 side, this ratio is approximately
$(Y-S_t)\big/\bigl(K_{d1}S_c/(S_c-Y)\bigr)$, which can be enormous because
$K_{d1}$ is small. Equivalently, bound \eqref{eq:x-bound} gives
$\rho'/\rho\le1+(A-A')/(Y_f+K_{d1})$, which is large when $A-A'$ exceeds
the free antibody $\sim\sqrt{K_{d1}S}$.

\textbf{Numerical illustration.} Take $K_{d1}=0.1$~nM, $K_{d2}=1000$~nM
($r=10^{-4}$), $A=90$~nM, $B=10$~nM, $M\ll K_{d2}$, and a treatment that
removes about 11\% of the target, $A'=80$~nM. Then $S_c=100$~nM and
$S_t=90$~nM, and $\sqrt{K_{d1}S}\approx3$~nM.
\begin{itemize}
    \item With $Y=95$~nM, $\delta_c=5$~nM and $\delta_t=-5$~nM. Substituting
    into \eqref{eq:x-yf},
    $$
    \begin{gathered}
    Y_f\approx\frac{-5+\sqrt{25+4(0.1)(100)}}{2}=\frac{-5+8.06}{2}\approx1.53~\text{nM},\\
    Y_f'\approx\frac{5+\sqrt{25+4(0.1)(90)}}{2}=\frac{5+7.81}{2}\approx6.41~\text{nM}.
    \end{gathered}
    $$
    By \eqref{eq:x-rho-eq}, $\rho\approx1.63/1000=1.6\times10^{-3}$ and
    $\rho'\approx6.51/1000=6.5\times10^{-3}$, so $\rho'/\rho\approx4$.
    (Solving \eqref{eq:x-F} numerically gives $\rho'/\rho\approx4.2$.)
    \item With $Y=50$~nM (deep in regime 1), the same 11\% drop gives
    $\rho'/\rho\approx\dfrac{90/40}{100/50}=1.125$.
\end{itemize}
An 11\% change in target abundance produces a fourfold change in the
relative off-target capture efficiency when the antibody is titrated near
equivalence, but only a 12.5\% change when the antibody is at half the
target concentration.

\textbf{Remark: an off-target sink.} If $M$ is not small compared with
$K_{d2}$ but off-target sites are still unsaturated ($Y_f\ll K_{d2}$), then
$M\theta_2\approx MY_f/K_{d2}$, and the antibody not on target is
$$
U=Y_f+M\theta_2\approx gY_f,
\qquad
g=1+\frac{M}{K_{d2}}.
$$
Repeating Step 1 with $U$ in place of $Y_f$ in \eqref{eq:r1-1} and
in antibody conservation, the free target is $\delta+U$ and the bound target
is $Y-U$, so mass action reads $K_{d1}(Y-U)=Y_f(\delta+U)$. Multiplying both
sides by $g$ and substituting $gY_f=U$,
$$
gK_{d1}(Y-U)=U(\delta+U).
$$
This is \eqref{eq:x-quad-exact} with $Y_f\to U$ and $K_{d1}\to gK_{d1}$.
All the results above carry over with these replacements; in particular,
$U\approx\sqrt{gK_{d1}S}$ at equivalence and $Y_f=U/g$. The off-target pool
acts as a buffer that soaks up a fixed fraction $1-1/g$ of the antibody not
on target, and it widens the window to $\sqrt{gK_{d1}S}$, which makes the
crossover less sharp.

\subsubsection{Off-target half-saturation crossover (between regimes 3 and 4)}

Once $Y-S\gg\sqrt{K_{d1}S}$, the target is saturated and the remaining
question is how the excess antibody distributes between solution and
off-target epitopes. Define the excess antibody
$$
E=Y-S>0.
$$

\textbf{Assumption.} The target is saturated: the free target
$S(1-\theta_1)=SK_{d1}/(Y_f+K_{d1})$ is negligible compared with $E$. In this
crossover $Y_f\sim K_{d2}\gg K_{d1}$, so the free target is about
$SK_{d1}/K_{d2}=rS$, and the assumption amounts to $rS\ll Y-S$.

\textbf{Step 1: a quadratic for $Y_f$.} Saturation means
$S\theta_1=S-S(1-\theta_1)\approx S$. Substituting this and
$\theta_2=Y_f/(Y_f+K_{d2})$ into antibody conservation,
$Y=Y_f+S+MY_f/(Y_f+K_{d2})$, and subtracting $S$ from both sides,
$$
E=Y_f+\frac{MY_f}{Y_f+K_{d2}}.
$$
In words, the excess antibody is either free or bound off-target.
Multiplying both sides by $Y_f+K_{d2}$,
$$
EY_f+EK_{d2}=Y_f^2+K_{d2}Y_f+MY_f,
$$
and collecting terms on one side,
\begin{equation}
Y_f^2+bY_f-EK_{d2}=0,
\qquad
b=K_{d2}+M-E.
\label{eq:x-quad-c}
\end{equation}

\textbf{Step 2: solve.} The product of the roots is $-EK_{d2}<0$, so there
is exactly one positive root. By the quadratic formula, and rationalizing
as in \eqref{eq:x-yf},
\begin{equation}
Y_f
=\frac{-b+\sqrt{b^2+4EK_{d2}}}{2}
=\frac{2EK_{d2}}{b+\sqrt{b^2+4EK_{d2}}},
\qquad
\rho=\frac{Y_f+K_{d1}}{Y_f+K_{d2}}\approx\frac{Y_f}{Y_f+K_{d2}},
\label{eq:x-yf-c}
\end{equation}
where the last approximation uses $K_{d1}\ll Y_f$.

\textbf{Step 3: the limits.}
\begin{itemize}
    \item \emph{Small excess, $E\ll K_{d2}+M$ (regime 3).} Then
    $b\approx K_{d2}+M>0$ and $4EK_{d2}\le4E(K_{d2}+M)\ll b^2$, so
    $\sqrt{b^2+4EK_{d2}}\approx b$. Substituting into the second form of
    \eqref{eq:x-yf-c},
    $$
    Y_f\approx\frac{2EK_{d2}}{2b}\approx\frac{K_{d2}(Y-S)}{K_{d2}+M}.
    $$
    This is $\ll K_{d2}$, so $\rho\approx(Y_f+K_{d1})/K_{d2}$; substituting
    $Y_f$ gives $\rho\approx r+(Y-S)/(K_{d2}+C+D)$, the regime-3 result.
    \item \emph{Large excess, $E\gg K_{d2}+M$ (regime 4).} Now
    $b=-(E-M-K_{d2})$ is large and negative. Factoring out $|b|$ and using
    the linear approximation $\sqrt{1+u}\approx1+u/2$ for small $u$,
    $$
    \sqrt{b^2+4EK_{d2}}
    =|b|\sqrt{1+\frac{4EK_{d2}}{b^2}}
    \approx|b|+\frac{2EK_{d2}}{|b|}
    \approx|b|+2K_{d2},
    $$
    where the last step uses $E/|b|\approx1$. Substituting into the first form
    of \eqref{eq:x-yf-c},
    $$
    Y_f\approx\frac{|b|+|b|+2K_{d2}}{2}=|b|+K_{d2}=E-M=Y-S-(C+D).
    $$
    All epitopes are saturated and the leftover antibody is free. Then
    $\rho\approx(E-M)/(E-M+K_{d2})\to1$ when $E-M\gg K_{d2}$, the regime-4
    result.
    \item \emph{Half-saturation.} $\rho\approx1/2$ when $Y_f=K_{d2}$, i.e.,
    when off-target epitopes are half occupied. Rather than solving
    \eqref{eq:x-quad-c}, substitute $Y_f=K_{d2}$ directly into
    $E=Y_f+MY_f/(Y_f+K_{d2})$:
    $$
    E_{1/2}=K_{d2}+\frac{M}{2},
    \qquad\text{i.e.,}\qquad
    Y=S+K_{d2}+\frac{C+D}{2}.
    $$
    The antibody must cover the target ($S$), half of the off-target pool
    ($M/2$), and leave $K_{d2}$ free.
\end{itemize}

\textbf{Two shapes of crossover.} The size of $M$ relative to $K_{d2}$
controls how the transition looks.
\begin{itemize}
    \item If $M\ll K_{d2}$, set $M=0$ in \eqref{eq:x-quad-c}:
    $Y_f^2+(K_{d2}-E)Y_f-EK_{d2}=0$. This factors as
    $(Y_f-E)(Y_f+K_{d2})=0$ (check by expanding:
    $Y_f^2+K_{d2}Y_f-EY_f-EK_{d2}$), so $Y_f=E$ and
    $$
    \rho\approx\frac{(Y-S)+K_{d1}}{(Y-S)+K_{d2}}.
    $$
    This is a gentle hyperbola: $\rho$ rises from $0.09$ to $0.91$ as
    $Y-S$ goes from $0.1K_{d2}$ to $10K_{d2}$, spanning two orders of
    magnitude.
    \item If $M\gg K_{d2}$, then \eqref{eq:x-quad-c} has exactly the form of
    \eqref{eq:x-quad-exact} with $Y_f$ unchanged, $K_{d1}\to K_{d2}$,
    $\delta\to M-E$, and $Y\to E$. So the antibody titrates the off-target
    pool the way it titrated the target in the equivalence crossover: there
    is a ``second equivalence'' at $E\approx M$ with window width
    $\sqrt{K_{d2}M}$. Half-saturation ($E=K_{d2}+M/2$) occurs before this
    second equivalence, and by $E=M$ we already have
    $Y_f\approx\sqrt{K_{d2}M}\gg K_{d2}$, so $\rho$ is close to 1.
\end{itemize}

\textbf{$\rho'$ versus $\rho$.} Because the target is saturated, $S$ enters
only through $E=Y-S$, and the treatment has
$E'=E+(A-A')>E$. Differentiating $E=Y_f+MY_f/(Y_f+K_{d2})$ implicitly with
respect to $E$,
$$
1=\frac{dY_f}{dE}\left(1+\frac{MK_{d2}}{(Y_f+K_{d2})^2}\right),
\qquad\text{so}\qquad
0<\frac{dY_f}{dE}\le1.
$$
Thus $Y_f'>Y_f$ and $\rho'>\rho$, as \eqref{eq:x-mono} requires. But the
effect is small. Near half-saturation, \eqref{eq:x-drho} with $Y_f=K_{d2}$
gives $d\rho/dY_f\approx K_{d2}/(2K_{d2})^2=1/(4K_{d2})$, so by the chain rule
$$
\rho'-\rho
\approx\frac{d\rho}{dY_f}\cdot\frac{dY_f}{dE}\cdot(A-A')
\le\frac{A-A'}{4K_{d2}},
\qquad
\frac{\rho'}{\rho}\lesssim1+\frac{A-A'}{2K_{d2}},
$$
using $\rho\approx1/2$ in the second expression. Unlike near equivalence,
$\rho$ is not small here, so a small absolute change in $\rho$ cannot become
a large relative change. With the parameters of the earlier illustration
and $C=D=1000$~nM, taking $Y=2500$~nM gives (numerically) $\rho\approx0.561$
and $\rho'\approx0.562$.

\subsubsection{Weak-binding crossover (between regimes 1 and 2)}

When $K_{d1}$ is comparable to $S$, neither ``all target bound'' nor
``almost no antibody bound'' holds. Clearing the denominators of
\eqref{eq:x-F} by multiplying both sides by $(Y_f+K_{d1})(Y_f+K_{d2})$,
$$
Y(Y_f+K_{d1})(Y_f+K_{d2})
=Y_f(Y_f+K_{d1})(Y_f+K_{d2})+SY_f(Y_f+K_{d2})+MY_f(Y_f+K_{d1}),
$$
which is a cubic in $Y_f$. Its formula is not illuminating, so in practice
one solves it numerically (Step 2 of the general principle guarantees a
unique positive solution). Two observations remove most of the need to do
so.

First, if off-target-bound antibody is negligible ($M\ll K_{d2}$), the
derivation of \eqref{eq:x-quad-exact} never used $K_{d1}\ll S$, so with
$\delta=S-Y$,
$$
Y_f=\frac{-(S-Y+K_{d1})+\sqrt{(S-Y+K_{d1})^2+4K_{d1}Y}}{2},
\qquad
\rho=\frac{Y_f+K_{d1}}{Y_f+K_{d2}}
$$
is exact across the whole crossover. (If $M$ is not negligible but
off-target sites are unsaturated, the remark on the off-target sink applies.)

Second, the regime-2 limit follows directly from conservation. Since
$\theta_1=Y_f/(Y_f+K_{d1})\le Y_f/K_{d1}$ and $Y_f\le Y$, the antibody bound
to target satisfies
$$
Y-Y_f\approx S\theta_1\le\frac{SY_f}{K_{d1}}\le\frac{S}{K_{d1}}\,Y,
$$
which is $\ll Y$ when $S\ll K_{d1}$, giving $Y_f\approx Y$. As $S/K_{d1}$
grows toward 1, a growing fraction of antibody is captured by the target,
$Y_f$ drops below $Y$, and $\rho$ drops from $(Y+K_{d1})/(Y+K_{d2})$ toward
regime 1's $r\,S/(S-Y)$ (if $Y<S$) or toward the equivalence or regime-3
behavior (if $Y\gtrsim S$). Throughout, \eqref{eq:x-mono} still guarantees
$\rho'>\rho$ when $A'<A$. The difference is small near regime 2: from the
implicit derivative in the general principle, $|dY_f/dS|\le\theta_1\le
Y_f/K_{d1}$, which is small when $Y_f\ll K_{d1}$.

\subsubsection{ChIP normalization}

From Regime 1's ChIP normalization, scaling by total spike-in reads
(ChIP-Rx) gives
$$
\frac{n_i'/N_s'}{n_i/N_s}
=\frac{a_i'}{a_i}\cdot\beta,
\qquad
\beta=\frac{B+\rho D}{B+\rho' D}.
$$
That derivation used only the read model, not any regime-specific
approximation for $\rho$, so it holds everywhere.

\textbf{The bias never favors the treatment.} By \eqref{eq:x-mono},
$\rho'>\rho$ whenever $A'<A$. Multiplying by $D\ge0$ and adding $B$,
$$
B+\rho' D\ge B+\rho D>0,
$$
so dividing both sides by $B+\rho'D$ gives $\beta\le1$, with equality only
if $D=0$ (or approximately, in the deep limits where $\rho'\approx\rho$). The
measured fold change is the true fold change multiplied by $\beta\le1$.
For a global decrease, $a_i'/a_i<1$, so the measured value is even smaller:
\textbf{ChIP-Rx with total spike-in reads never understates, and generally
overstates, a global decrease}, in every regime. (Equivalently, every locus
in the treatment is biased downward by the same factor $\beta$.)

\textbf{How large can the bias be?} Compare $\beta$ with $\rho/\rho'$.
Putting both over a common denominator,
$$
\beta-\frac{\rho}{\rho'}
=\frac{\rho'(B+\rho D)-\rho(B+\rho' D)}{\rho'(B+\rho' D)}
=\frac{\rho' B+\rho\rho' D-\rho B-\rho\rho' D}{\rho'(B+\rho' D)}
=\frac{B(\rho'-\rho)}{\rho'(B+\rho' D)}
\ge0.
$$
Therefore
$$
\frac{\rho}{\rho'}\le\beta\le1.
$$
The bias approaches its worst case $\rho/\rho'$ when $\rho'D\gg B$, i.e.,
when most spike-in reads come from off-target spike-in epitopes, and
vanishes when $\rho'D\ll B$. Combined with the results above:
\begin{itemize}
    \item Deep in regimes 1 (with $Y\ll S$), 2, and 4, $\rho'\approx\rho$,
    so $\beta\approx1$.
    \item Near equivalence, $\rho'/\rho$ can be large (fourfold in the
    illustration), so $\beta$ can be far below 1 if $\rho'D$ is comparable
    to or larger than $B$. This is where the bias is largest.
    \item Near off-target half-saturation, $\rho'/\rho\lesssim1+(A-A')/(2K_{d2})$,
    so $\beta$ is close to 1 unless $A-A'$ is comparable to $K_{d2}$.
\end{itemize}

\textbf{Numerical illustration.} Use $K_{d1}=0.1$~nM, $K_{d2}=1000$~nM,
$A=90$~nM, $A'=80$~nM, $B=10$~nM, and now a substantial off-target pool
$C=D=1000$~nM (so $g=3$ in the off-target-sink remark). The true global fold
change is $80/90\approx0.89$. Solving \eqref{eq:x-F} numerically:
\begin{itemize}
    \item $Y=50$~nM (regime 1): $\beta\approx0.998$.
    \item $Y=95$~nM (near equivalence): $\rho\approx1.20\times10^{-3}$,
    $\rho'\approx2.82\times10^{-3}$, so $\rho D\approx1.2$~nM and
    $\rho'D\approx2.8$~nM, giving $\beta\approx11.2/12.8\approx0.87$. The
    measured fold change is $0.89\times0.87\approx0.78$. Scanning $Y$, the
    smallest $\beta$ ($\approx0.83$) occurs at $Y\approx107$~nM, just past
    equivalence (the off-target sink shifts it slightly).
    \item $Y=2500$~nM (off-target half-saturation): $\beta\approx0.998$.
\end{itemize}
As in Regime 1, counting only reads at the spike-in's target loci
($N_{s,B}$) removes the dependence on $\rho$ and $\rho'$ entirely, and
therefore removes this bias in every regime.
